\section{Deliverables of the Project}

A deliverable is a concrete output that can be reviewed, assessed, and checked off \cite{del1}. Vague outputs like ``work in progress'' or ``mostly done'' do not count. Every deliverable listed here has a clear definition of done. The five deliverable categories below cover the full spectrum of FYP outputs.

\subsection{Requirements and Analysis Deliverables}

The early deliverables from this project include the problem analysis, the SRS document, stakeholder interview notes, use-case descriptions, and the literature review of existing systems. These outputs establish the academic justification for the project. They explain why the current state of lab exam management is inadequate and what the proposed system will do differently. Without these, later design choices have no documented rationale.

\subsection{Design Deliverables}

Design deliverables translate requirements into buildable plans. For SmartExam, this means system architecture diagrams, ER diagrams, class diagrams, use-case diagrams, data-flow diagrams, and the scheduling artifacts (Gantt chart, WBS, CPM network). These outputs reduce the gap between ``what we need'' and ``how we build it''.

\subsection{Prototype and Module Deliverables}

The primary software deliverable is a working prototype of all twelve modules. Specifically:

\begin{itemize}[noitemsep]
\item A WPF desktop client for students (HWID binding, fullscreen lockdown, process telemetry, auto-save)
\item A React single-page application for teachers and administrators (exam management, live monitoring grid, reporting)
\item An ASP.NET Core REST API server connected to a Neon PostgreSQL database (authentication, scheduling, submission handling, audit logging)
\end{itemize}

Each module must run end-to-end before it is counted as delivered. A UI screen without a working backend endpoint does not qualify.

\subsection{Testing and Validation Deliverables}

Testing deliverables include module-level test results, integration test reports, and user flow validation notes. These outputs show that the prototype has been examined critically, not just assembled. They also document which features were fully tested, which were partially tested, and which remain open to future refinement.

\subsection{Documentation and Presentation Deliverables}

The final FYP report, the academic presentation deck, and the supervisor review responses are all deliverables. These are not afterthoughts. A strong system that is poorly explained will score below its potential. Documentation must be completed incrementally — one chapter per sprint cycle — so that writing never falls more than one sprint behind implementation.
